\documentclass[a4paper, 14pt]{extarticle}

\usepackage{cmap}
\usepackage[T2A]{fontenc}
\usepackage[utf8]{inputenc}
\usepackage[english, russian]{babel}
\usepackage{amsmath}
\usepackage{subcaption}
\usepackage{graphicx}
\usepackage{float}
\usepackage{hyperref}
\usepackage{multirow}

\usepackage[left=3cm, right=1.5cm, top=2cm, bottom=2cm]{geometry}
\usepackage{setspace}
\onehalfspacing 

\usepackage{indentfirst}
\setlength{\parindent}{1.25cm} 
\setlength{\emergencystretch}{3em}
\usepackage[expansion=false,protrusion=false,nopatch=footnote]{microtype}

\usepackage{tocloft}
\addto\captionsrussian{\renewcommand{\contentsname}{СОДЕРЖАНИЕ}}
\renewcommand{\cfttoctitlefont}{\hfill\normalfont}
\renewcommand{\cftaftertoctitle}{\hfill\vspace{\baselineskip}}
\renewcommand{\cftsecdotsep}{\cftdotsep}

\begin{document}

\begin{titlepage}
    \begin{center}
        Министерство науки и высшего образования Российской Федерации\\
        Федеральное государственное автономное образовательное учреждение высшего образования\\
        «Санкт-Петербургский политехнический университет Петра Великого»\\
        
        \textbf{Физико-механический институт}\\
        \textbf{Высшая школа прикладной математики и вычислительной физики}
    \end{center}

    \vfill
    
    \begin{center}
        \textbf{Отчёт по лабораторной работе \textnumero~ 8\\[0.5cm]}
        \textbf{Доверительные интервалы. Критерий Стьюдента и F-тест\\}
        по дисциплине <<Математическая статистика>>
    \end{center}

    \vfill

    \hfill
    \begin{minipage}{0.45\textwidth}
        \textbf{Выполнил студент гр. 5030102/30202:}\\
        Шильдт Д.С.\\[0.5cm]
        \textbf{Преподаватель:}\\ 
        Баженов А.Н.
    \end{minipage}

    \vfill

    \begin{center}
        Санкт-Петербург\\2026
    \end{center}
\end{titlepage}

\setcounter{page}{2}

\tableofcontents
\newpage

\section{Постановка задачи}

\textbf{Задание:}
\begin{enumerate}
    \item Сгенерировать две выборки объёмом $n_1 = 20$ и $n_2 = 100$ из $N(0,1)$.
    \item Построить доверительные интервалы для:
    \begin{itemize}
        \item математического ожидания (на основе $\bar{x}$ и $t$-распределения Стьюдента)
        \item дисперсии на основе ($s^2$ и $\chi^2$-распределения).
    \end{itemize}
    \item Проверить гипотезу о равенстве дисперсий двух независимых выборок с помощью $F$-теста (критерий Фишера).
\end{enumerate}

\textbf{Цели и задачи:}
\begin{itemize}
    \item Освоить построение интервальных оценок параметров.
    \item Исследовать связь между точечными и интервальными оценками.
\end{itemize}

\section{Теоретическая часть}

\subsection{Доверительные интервалы для параметров нормального распределения}

Пусть $x_1, x_2, \dots, x_n$ --- выборка из нормальной генеральной совокупности $N(m, \sigma^2)$. В качестве точечных оценок параметров используются выборочное среднее:
\begin{equation*}
    \bar{x} = \frac{1}{n} \sum_{i=1}^{n} x_i,
\end{equation*}
и выборочная дисперсия:
\begin{equation*}
    s^2 = \frac{1}{n} \sum_{i=1}^{n} (x_i - \bar{x})^2.
\end{equation*}

Для математического ожидания применяется \textit{статистика Стьюдента:}
\begin{equation*}
    T=\sqrt{n-1}\cdot\frac{\overline{x}-m}{s},
\end{equation*}
которая при справедливости предположения о нормальности распределена по закону Стьюдента с $n-1$ степенями свободы. 

Доверительный интервал для $m$ с доверительной вероятностью $\gamma = 1 - \alpha$:
\begin{equation*}
    P\left(\overline{x}-\frac{st_{1-\alpha/2}(n-1)}{\sqrt{n-1}}<m<\overline{x}+\frac{st_{1-\alpha/2}(n-1)}{\sqrt{n-1}}\right)=1-\alpha,
\end{equation*}
где $t_{1-\alpha/2}(n-1)$ --- \textit{квантиль распределения Стьюдента} c $n-1$ степенями свободы и порядка $1-\alpha/2$. 

Построение доверительного интервала для неизвестной дисперсии $\sigma^2$ основано на свойствах распределения хи-квадрат. Известно, что для выборки объема $n$ из нормальной генеральной совокупности случайная величина $\frac{n s^2}{\sigma^2}$ подчиняется закону $\chi^2$ с $n-1$ степенями свободы:$$\frac{n s^2}{\sigma^2} \sim \chi^2(n-1)$$ Одна степень свободы в данном случае теряется из-за того, что при расчете выборочной дисперсии $s^2$ вместо истинного математического ожидания используется его оценка — выборочное среднее $\bar{x}$. Для заданной доверительной вероятности $\gamma = 1 - \alpha$ доверительный интервал:
\begin{equation*}
    P\left(\frac{s\sqrt{n}}{\sqrt{\chi_{1-\alpha/2}^2(n-1)}}<\sigma<\frac{s\sqrt{n}}{\sqrt{\chi_{\alpha/2}^2(n-1)}}\right)=1-\alpha.
\end{equation*}

\subsection{Критерий Фишера}

Для проверки гипотезы $H_0: \sigma_1^2 = \sigma_2^2$ о равенстве дисперсий двух независимых нормальных выборок используется $F$-критерий. В качестве статистики берется отношение исправленных выборочных дисперсий:
\begin{equation*}
    s_1^2 = \frac{1}{n_1 - 1} \sum_{i=1}^{n_1} (x_i - \bar{x})^2, \qquad
    s_2^2 = \frac{1}{n_2 - 1} \sum_{i=1}^{n_2} (y_i - \bar{y})^2.
\end{equation*}
Тогда
\begin{equation*}
    F = \frac{s_1^2}{s_2^2},
\end{equation*}
где при справедливости $H_0$ величина $F$ имеет распределение Фишера с числом степеней свободы $n_1-1$ и $n_2-1$. Для двусторонней проверки гипотеза отвергается, если
\begin{equation*}
    F < F_{\alpha/2}(n_1-1, n_2-1) \quad \text{или} \quad F > F_{1-\alpha/2}(n_1-1, n_2-1).
\end{equation*}

\section{Реализация}

\subsection{Язык программирования и среда}
\begin{itemize}
    \item \textbf{Используемый язык:} Python 3.14.2
    \item \textbf{Используемая среда разработки:} Visual Studio Code
\end{itemize}

\subsection{Используемые библиотеки}
\begin{itemize}
    \item \textbf{numpy} --- используется для генерации нормальных выборок и вычисления выборочных характеристик.
    \item \textbf{scipy.stats} --- применяется для вычисления квантилей распределений Стьюдента, $\chi^2$ и $F$.
    \item \textbf{pathlib} --- обеспечивает формирование путей к каталогам данных.
    \item \textbf{csv} --- используется для сохранения результатов вычислений в табличном виде.
\end{itemize}

\subsection{Суть реализованных методов}

\begin{itemize}
    \item \texttt{generate\_normal\_sample(n, rng)} --- формирует выборку из распределения $N(0,1)$ заданного объёма.
    \item \texttt{mean\_confidence\_interval(sample, alpha)} ---\\ вычисляет доверительный интервал для математического ожидания по $t$-распределению.
    \item \texttt{variance\_confidence\_interval(sample, alpha)} ---\\ строит доверительный интервал для дисперсии по $\chi^2$-распределению.
    \item \texttt{fisher\_variance\_test(sample\_1, sample\_2, alpha)} ---\\ реализует двусторонний $F$-тест для проверки равенства дисперсий.
    \item \texttt{save\_samples(...)} и \texttt{save\_interval\_summary(...)} ---\\ сохраняют исходные данные и итоговые результаты в CSV-файлы.
\end{itemize}

\section{Результаты}

Для вычислений использовалось значение уровня значимости $\alpha = 0.05$ и фиксированное зерно генератора случайных чисел $42$. Ниже приведены выборочные оценки и соответствующие доверительные интервалы.

\begin{table}[H]
    \centering
    \caption{Доверительные интервалы для параметров двух выборок из $N(0,1)$}
    \label{tab:ci_summary}
    \resizebox{\textwidth}{!}{%
    \begin{tabular}{|c|c|c|c|c|c|}
        \hline
        \textbf{Выборка} & $\mathbf{n}$ & $\mathbf{\bar{x}}$ & \textbf{ДИ для $\mathbf{\mu}$} & $\mathbf{s^2}$ & \textbf{ДИ для $\mathbf{\sigma^2}$} \\ \hline
        $n_1$ & 20 & $-0.0329$ & $[-0.4402;\,0.3743]$ & $0.7193$ & $[0.4379;\,1.6153]$ \\ \hline
        $n_2$ & 100 & $-0.0659$ & $[-0.2190;\,0.0872]$ & $0.5890$ & $[0.4587;\,0.8029]$ \\ \hline
    \end{tabular}%
    }
\end{table}

В таблице \ref{tab:fisher_summary} приведён результат проверки гипотезы о равенстве дисперсий двух независимых выборок по $F$-критерию.

\begin{table}[H]
    \centering
    \caption{Результат $F$-теста для двух выборок}
    \label{tab:fisher_summary}
    \small
    \setlength{\tabcolsep}{8pt}
    \renewcommand{\arraystretch}{1.2}
    \begin{tabular}{|c|c|c|c|c|}
        \hline
        \shortstack{\textbf{$F$-статистика}\\\textbf{$s_1^2/s_2^2$}} &
        \shortstack{\textbf{$df_1$}} &
        \shortstack{\textbf{$df_2$}} &
        \textbf{Критический интервал} &
        \shortstack{\textbf{p-value}\\\textbf{Решение}} \\ \hline
        $1.2726$ & 19 & 99 & $(0.4510;\,1.8696)$ & $0.4381$\\
        \hline
    \end{tabular}
    
    \vspace{0.4cm}
    \begin{minipage}{0.92\textwidth}
        \centering
        \textit{Вывод:} гипотеза $H_0$ о равенстве дисперсий не отвергается.
    \end{minipage}
\end{table}

Ширина доверительного интервала для $m$ заметно уменьшается при увеличении объёма выборки с $20$ до $100$, что соответствует свойствам оценки по распределению Стьюдента. Проверка гипотезы о равенстве дисперсий показала, что отношение исправленных выборочных дисперсий не выходит за пределы критической области, поэтому на уровне значимости $\alpha = 0.05$ гипотеза $H_0$ не отвергается.

\section{Обсуждение}

Результаты работы подтверждают, что точечные и интервальные оценки связаны между собой: одно и то же выборочное среднее $\bar{x}$ может соответствовать узкому или широкому доверительному интервалу в зависимости от объёма выборки и разброса наблюдений. В случае дисперсии влияние случайных колебаний проявляется сильнее, что хорошо видно по второй выборке: интервал для $\sigma^2$ уже не включает единицу, хотя исходное распределение задавалось как $N(0,1)$.

$F$-критерий в данной постановке оказывается корректным инструментом для сравнения двух дисперсий, поскольку обе выборки получены из нормального закона. При этом сама проверка чувствительна к форме критической области и требует соблюдения двусторонней схемы принятия решения.

\section{Выводы}

В ходе выполнения работы были построены доверительные интервалы для математического ожидания и дисперсии двух независимых выборок из $N(0,1)$ объёмом $20$ и $100$ наблюдений. Для обоих наборов данные интервалов подтвердили близость выборочных оценок к теоретическому значению $\mu = 0$.

Проверка гипотезы о равенстве дисперсий с помощью $F$-теста показала, что статистика $F = 1.2726$ не выходит за пределы критической области, поэтому гипотеза $H_0$ не отвергается. Это означает, что различие между выборочными дисперсиями статистически незначимо на уровне $\alpha = 0.05$.

\refstepcounter{section}
\addcontentsline{toc}{section}{\thesection\ Список литературы}
\renewcommand{\refname}{\thesection\ Список литературы}

\begin{thebibliography}{9}

    \bibitem{babakhina}
    Бабахина С. А., Баженов А. Н. Практикум по математической статистике: учеб. пособие. — СПб. : Санкт-Петербургский политехнический университет Петра Великого, 2026. 

    \bibitem{maksimov}
    Максимов Ю. Д. Математика. Теория и практика по математической статистике. Конспект-справочник по теории вероятностей : учеб. пособие. — СПб. : Изд-во Политехн. ун-та, 2009. — 395 с. 

\end{thebibliography}

\section{Приложение}
Исходный код программ доступен в репозитории на GitHub: \\
\url{https://github.com/Reichenau/math-stats-labs}
 
\end{document}
